\section{Per-robot onboarding detail}
\label{app:onboarding}

Table~\ref{tab:onboarding-detail} reports the per-robot
breakdown behind the $7.8 \pm 1.6$ minute aggregate in
\S\ref{sec:onboarding}; the $\$2.98 \pm \$0.84$ cost figure is
recorded per run rather than per robot, so it appears in
aggregate only.
The robot lineup is in main paper Figure~\ref{fig:hero} (top).

\noindent\textbf{Behavioral-test strictness.} The Validate
harness is deterministic framework code, not an LLM call
(Appendix~\ref{app:prompts}); however, the harness version used
for these headline onboarding runs graded the quadruped
behavioral tests leniently: \texttt{walk\_forward} was scored
ran-without-exception rather than by displacement, and Go2's
\texttt{sit} passes despite leaving body height at 0.28\,m
(A1 and ANYmal-C genuinely lower to 0.11\,m and 0.09\,m). The
released harness grades both on physics (\texttt{sit} must drop
below 0.8$\times$ standing height, \texttt{walk\_forward} must
displace the base $>$3\,cm while upright); we report the runs
under the harness they were validated with. The structural tests
(no-skeleton import, build, \texttt{home}) are unaffected, and
every downstream task number in \S\ref{sec:capability} is graded
by the benchmark's independent physics harness, never by the
onboarding harness; the Go2 smoke result in
Appendix~\ref{app:smoke} (\texttt{stand\_up} 0/2) shows the
leniency does not propagate to task grading.

\noindent\textbf{Strict re-validation.} Re-running the released
physics-graded harness on the same eight drivers quantifies the
gap. All eight keep every structural test. The five arms pass all
behavioral tests but fail the scene-perception API check, a
requirement added to the harness after these drivers were
synthesized (the \S\ref{sec:cap-portability} synthesis runs do
require and pass it). The three quadrupeds each fail one to two
behavioral tests: Go2 \texttt{sit} leaves body height at 0.28\,m
(threshold 0.22\,m), A1 \texttt{walk\_forward} moves $-2.7$\,cm,
and ANYmal-C drops from 0.61\,m to 0.47\,m during
\texttt{stand\_up} and does not displace on
\texttt{walk\_forward}. The headline 8/8 therefore reads as
structural plus smoke-level behavioral validation, exactly as
stated; under the released harness it is 0/8 fully passing: the
arms miss only the later-added scene-perception check, and the
quadrupeds fail locomotion.

Table~\ref{tab:provenance} maps every headline result block in the paper
to the driver it ran on, the delivery mode (\S\ref{sec:modes}),
and any human code involved after synthesis, so that no headline
number has ambiguous provenance.

\begin{table*}[t]
\centering
\footnotesize
\setlength{\tabcolsep}{3pt}
\begin{tabular}{llll}
\toprule
\textbf{Result block} & \textbf{Robot(s)} & \textbf{Driver provenance} & \textbf{Post-synthesis code} \\
\midrule
\S\ref{sec:onboarding} onboarding scoreboard & 8 robots & self-contained, one per robot (Sonnet~4.5) & none \\
\S\ref{sec:cap-cross} cross-skill, \S\ref{sec:cap-rlsteelman} multi-task & SO-101 & spec mode (skeleton + agent spec, Sonnet~4.5) & none \\
\S\ref{sec:cap-portability} ``uses API'' (A-A, CaP, CaP\textsubscript{fs}) & SO-101 & the \S\ref{sec:cap-cross} spec-mode driver, frozen & none \\
\S\ref{sec:cap-portability} ``writes API'' / Self & SO-101 & self-contained, one per evaluated model & none \\
App.\,\ref{app:openvla-sweep} OpenVLA & SO-101 & n/a (policy); translation via the frozen driver & n/a \\
\S\ref{sec:cap-cross-robot} simple suite & Piper, UR5e & self-contained (Sonnet~4.5) & none \\
\S\ref{sec:cap-cross-robot} end-to-end Pick & Piper & self-contained vs.\ pickbench scene & 13-line patch (\S\ref{sec:cap-cross-robot}) \\
\S\ref{sec:cap-cross-robot} locomotion v2 & Go2, ANYmal-C & shared self-contained (Sonnet~4.5) & none \\
\S\ref{sec:scoping-franka} reach drill-down & Franka & self-contained (Sonnet~4.5) & none \\
Appendix~\ref{app:morphology} aerial, humanoid & Skydio~X2, H1 & self-contained, class-specific prompt & none \\
\bottomrule
\end{tabular}
\caption{Driver provenance for every result block. ``Spec mode''
= the agent writes a JSON spec consumed by a 600-line skeleton
template; ``self-contained'' = the agent writes the entire driver
file (\S\ref{sec:modes}). The frozen spec-mode reference driver
is what all seven models operate in the
\S\ref{sec:cap-portability} ``uses API'' comparison.}
\label{tab:provenance}
\end{table*}

\begin{table*}[t]
\centering
\small
\begin{tabular}{lrrrrr}
\toprule
\textbf{Robot} & \textbf{Joints} & \textbf{Class}
   & \textbf{Driver LOC} & \textbf{Validate tests}
   & \textbf{Wall (min)} \\
\midrule
SO-101       &  5 & arm + gripper   & 546 & 7/7 &  5.1 \\
Piper        &  6 & arm + gripper   & 377 & 7/7 &  6.4 \\
UR5e         &  6 & arm             & 423 & 7/7 &  6.8 \\
Franka Panda &  7 & arm + gripper   & 527 & 7/7 &  8.4 \\
KUKA iiwa14  &  7 & arm             & 510 & 7/7 &  8.0 \\
Unitree Go2  & 12 & quadruped       & 546 & 8/8 &  9.1 \\
Unitree A1   & 12 & quadruped       & 458 & 8/8 &  8.6 \\
ANYmal-C     & 12 & quadruped       & 427 & 8/8 & 10.0 \\
\midrule
\textbf{Mean} &    &                & 477 & all pass & \textbf{7.8} \\
\bottomrule
\end{tabular}
\caption{Per-robot synthesis and validation. LOC is the
agent-written \texttt{driver\_from\_scratch.py}; arms run a 7-test
harness, quadrupeds an 8-test harness (locomotion in place of IK and
gripper). Mean wall-clock matches the $7.8 \pm 1.6$\,min headline
(\S\ref{sec:onboarding}).}
\label{tab:onboarding-detail}
\end{table*}

